\chapter{Contraintes Temporelles, Fréquences \& Analyse Temps Réel}
\label{chap:realtime_timings}

\section{Tableau Récapitulatif des Fréquences et Périodicités Système}
La conception en microservices asynchrones impose un cadencement rigoureux de chaque boucle d'exécution afin de concilier la réactivité face aux événements imprévus et la bande passante limitée des bus matériels. 

Le Tableau~\ref{tab:frequencies_master} synthétise les fréquences nominales, les périodes associées, les mécanismes de synchronisation et les chemins critiques de chaque composant de la pile FastDrink.

\begin{table}[htbp]
\centering
\footnotesize
\begin{tabularx}{\textwidth}{lccccX}
\toprule
\textbf{Composant / Service} & \textbf{Fréquence} & \textbf{Période ($T$)} & \textbf{Gigue Max.} & \textbf{Mécanisme} & \textbf{Chemin Critique / Rôle} \\
\midrule
\textbf{\texttt{ServiceManager}} & 10~Hz & 100~ms & $\pm 5$~ms & \code{usleep()} & Supervision démons, lien Ethernet filaire, Dropbear SSH. \\
\textbf{\texttt{FSM}} (Contrôle) & \textbf{50~Hz} & \textbf{20~ms} & $\pm 0{,}5$~ms & \code{sleep\_until} & Décision, régulateur PD, profileur de vitesse trapézoïdal. \\
\textbf{\texttt{Motors}} (Loop) & 100~kHz & 10~$\mu$s & $\pm 2$~$\mu$s & \code{sleep\_until} & Polling SHM, slew-rate limiter de direction continu. \\
\textbf{PCA9685 PWM} & 50~Hz & 20~ms & Matériel & Hardware & Période PWM servomoteur RC (impulsions 1 à 2~ms). \\
\textbf{\texttt{Sensors}} (Loop) & 100~kHz & 10~$\mu$s & $\pm 2$~$\mu$s & \code{sleep\_until} & Scrutation I2C ligne, calcul barycentrique, mise à jour SHM. \\
\textbf{HC-SR04 Sonar} & \textbf{50~Hz} & \textbf{20~ms} & $\pm 1$~ms & Diviseur $N$ & Impulsion 10~$\mu$s, écoute écho (5 à 25~ms max). \\
\textbf{\texttt{PeripheralsService}}& 200~Hz & 5~ms & $\pm 0{,}5$~ms & \code{poll(5ms)} & Échantillonnage boutons (anti-rebond $\ge 10$~ms), buzzer. \\
\textbf{\texttt{TelemetryService}} & 50~Hz & 20~ms & $\pm 1$~ms & \code{sleep\_until} & Collecte SHM, dérivation d'accélération, UDP broadcast. \\
\textbf{Console Télémétrie} & 5~Hz & 200~ms & $\pm 10$~ms & Modulo loop & Affichage synthétique des métriques sur terminal console. \\
\textbf{\texttt{SimBridge} (Godot)} & 50~Hz & 20~ms & $\pm 0{,}5$~ms & Socket TCP & Synchronisation avec \code{\_physics\_process} de Godot 4. \\
\bottomrule
\end{tabularx}
\caption{Table maîtresse des fréquences d'exécution et contraintes temporelles.}
\label{tab:frequencies_master}
\end{table}

\section{Stratégie d'Isolation Temporelle des Périphériques Lents}
L'un des défis majeurs dans la programmation temps réel sous Linux consiste à éviter qu'une opération d'entrée-sortie lente ne bloque une boucle de contrôle rapide.

Le capteur ultrasonique HC-SR04 en est l'illustration parfaite~: la mesure par écho sonore requiert entre $1\text{ ms}$ (pour un obstacle à 17~cm) et $25\text{ ms}$ (pour la portée maximale). Si cette lecture était exécutée directement à l'intérieur de la boucle rapide de \code{Sensors}, le thread subirait un gel imprévisible, rendant impossible la détection haute fréquence de la ligne optique.

La solution mise en œuvre repose sur une \textbf{isolation de fréquence synchrone par diviseur d'horloge}~:
\begin{equation}
    N = \left\lfloor \frac{f_{\text{loop}}}{f_{\text{ultrason}}} \right\rfloor = \left\lfloor \frac{100\,000}{50} \right\rfloor = 2\,000 \text{ cycles}
\end{equation}
La mesure acoustique n'est déclenchée que lorsque $\text{loop\_tick} \pmod N = 0$. Durant les 1999 autres cycles, le service réutilise la dernière distance filtrée tout en continuant d'échantillonner le ruban noir à pleine vitesse sans aucun retard.

\section{Chronogramme d'Exécution Temps Réel}
La Figure~\ref{fig:timing_diagram} illustre le cadencement relatif des processus sur un intervalle type de 40~ms (deux trames de contrôle à 50~Hz).

\begin{figure}[htbp]
\centering
\begin{tikzpicture}[
    scale=0.78, every node/.style={transform shape},
    timeline/.style={->, thick, >=stealth},
    gridline/.style={help lines, dashed, gray!50},
    slot/.style={draw=black!70, fill=blue!15, rounded corners=2pt, minimum height=0.45cm, font=\scriptsize},
    sonarslot/.style={draw=green!60!black, fill=green!20, rounded corners=2pt, minimum height=0.45cm, font=\scriptsize},
    fsmslot/.style={draw=purple!70!black, fill=purple!20, rounded corners=2pt, minimum height=0.45cm, font=\scriptsize},
    motorslot/.style={draw=orange!70!black, fill=orange!20, rounded corners=2pt, minimum height=0.45cm, font=\scriptsize}
]

% Axes temporels
\draw[timeline] (0, 4) -- (13.5, 4) node[right, font=\footnotesize] {$t$ (ms)};
\draw[timeline] (0, 2.7) -- (13.5, 2.7) node[right, font=\footnotesize] {$t$ (ms)};
\draw[timeline] (0, 1.4) -- (13.5, 1.4) node[right, font=\footnotesize] {$t$ (ms)};
\draw[timeline] (0, 0.1) -- (13.5, 0.1) node[right, font=\footnotesize] {$t$ (ms)};

% Labels de canaux
\node[anchor=east, font=\small\bfseries] at (-0.2, 4.3) {Sensors (I2C Ligne)};
\node[anchor=east, font=\small\bfseries] at (-0.2, 3.0) {Sensors (HC-SR04 Sonar)};
\node[anchor=east, font=\small\bfseries] at (-0.2, 1.7) {FSM (Régulation 50~Hz)};
\node[anchor=east, font=\small\bfseries] at (-0.2, 0.4) {Motors (PCA9685 PWM)};

% Lignes de repère temporel
\foreach \x/\t in {0/0, 3/10, 6/20, 9/30, 12/40} {
    \draw[gridline] (\x, -0.3) -- (\x, 4.8);
    \node[above, font=\tiny\bfseries] at (\x, 4.8) {\t~ms};
}

% Activités Sensors I2C (haute fréquence régulière)
\foreach \x in {0.2, 0.8, 1.4, 2.0, 2.6, 3.2, 3.8, 4.4, 5.0, 5.6, 6.2, 6.8, 7.4, 8.0, 8.6, 9.2, 9.8, 10.4, 11.0, 11.6, 12.2} {
    \draw[fill=blue!30, draw=blue!80] (\x, 3.9) rectangle ++(0.35, 0.4);
}

% Activité HC-SR04 Sonar (rafale 50 Hz avec temps de vol)
\draw[sonarslot, minimum width=1.5cm] (0.8, 2.95) rectangle ++(1.4, 0.45) node[midway] {Echo (7~ms)};
\draw[sonarslot, minimum width=1.3cm] (6.8, 2.95) rectangle ++(1.2, 0.45) node[midway] {Echo (6~ms)};

% Activité FSM (déclenchée à chaque frontière 20 ms)
\draw[fsmslot, minimum width=1.6cm] (2.3, 1.65) rectangle ++(1.8, 0.45) node[midway] {PD + Profiler};
\draw[fsmslot, minimum width=1.6cm] (8.3, 1.65) rectangle ++(1.8, 0.45) node[midway] {PD + Profiler};

% Activité Motors PCA9685 (émission PWM hardware 50 Hz)
\draw[motorslot, minimum width=2.4cm] (0.1, 0.35) rectangle ++(2.0, 0.45) node[midway] {PWM Période 1 (20~ms)};
\draw[motorslot, minimum width=2.4cm] (6.1, 0.35) rectangle ++(2.0, 0.45) node[midway] {PWM Période 2 (20~ms)};

% Flèches de synchronisation IPC
\draw[-{Stealth}, thick, purple!80!black] (2.2, 3.8) to[out=-90, in=90] (2.5, 2.1);
\node[right, font=\tiny, purple!80!black] at (2.4, 2.8) {SHM Write};

\draw[-{Stealth}, thick, orange!80!black] (4.1, 1.65) to[out=-90, in=90] (4.5, 0.8);
\node[right, font=\tiny, orange!80!black] at (4.5, 1.2) {Consigne $\mathbf{u}$};

\end{tikzpicture}
\caption{Chronogramme temps réel illustrant l'isolation de fréquence et l'absence d'inversion de priorité.}
\label{fig:timing_diagram}
\end{figure}

\section{Justification Mathématique des Limites Physiques FastDrink}

\subsection{\texorpdfstring{Démonstration du Seuil Critique d'Accélération ($a_{\max} \le 0{,}735\text{ m/s}^2$)}{Démonstration du Seuil Critique d'Accélération (a <= 0.735 m/s2)}}
Considérons une bille de masse $m$ reposant au fond d'une coupelle sphérique de profondeur utile $\Delta h = 1{,}50\text{ mm} = 0{,}0015\text{ m}$ et de rayon horizontal bord-à-centre $d = 20{,}0\text{ mm} = 0{,}020\text{ m}$.

Le travail mécanique fourni par la force d'inertie lors d'une accélération linéaire $a$ sur la distance $d$ pour franchir la bordure s'écrit~:
\begin{equation}
    W_{\text{inertie}} = F_{\text{inertie}} \cdot d = (m \cdot a) \cdot d
\end{equation}
L'énergie potentielle de pesanteur nécessaire pour faire monter la bille sur la hauteur $\Delta h$ est~:
\begin{equation}
    \Delta E_p = m \cdot g \cdot \Delta h
\end{equation}
À la limite du basculement, par le théorème de l'énergie mécanique ($W_{\text{inertie}} = \Delta E_p$)~:
\begin{equation}
    m \cdot a \cdot d = m \cdot g \cdot \Delta h \implies a = g \cdot \frac{\Delta h}{d}
\end{equation}
En retenant l'accélération gravitationnelle standard $g = 9{,}80665\text{ m/s}^2$~:
\begin{equation}
    a_{\max} = 9{,}80665 \times \frac{0{,}0015}{0{,}020} = 0{,}735498\text{ m/s}^2 \approx \mathbf{0{,}735\text{ m/s}^2}
\end{equation}
Toute accélération instantanée supérieure à cette valeur éjecte irrémédiablement la bille.

\subsection{\texorpdfstring{Quantum Maximal de Vitesse par Pas Temporel ($\Delta t = 20\text{ ms}$)}{Quantum Maximal de Vitesse par Pas Temporel (Delta t = 20 ms)}}
Pour la boucle décisionnelle cadencée à 50~Hz ($\Delta t = 0{,}020\text{ s}$), le saut maximal de vitesse permis par le \code{MotionProfiler} entre deux itérations consécutives est rigoureusement borné par~:
\begin{equation}
    \Delta v_{\text{cycle}} \le a_{\max} \cdot \Delta t = 0{,}735 \times 0{,}020 = 0{,}0147\text{ m/s} \quad (1{,}47\text{ cm/s par cycle})
\end{equation}

\subsection{\texorpdfstring{Vitesse Maximale Admissible en Courbe ($R_{\min} = 120\text{ mm}$)}{Vitesse Maximale Admissible en Courbe (Rmin = 120 mm)}}
Lors du franchissement d'un virage de rayon minimal imposé $R_{\min} = 120\text{ mm} = 0{,}12\text{ m}$ (Spécification S3), l'accélération centripète $a_c = \frac{v^2}{R}$ ne doit pas excéder $a_{\max}$~:
\begin{equation}
    v_{\text{virage}} \le \sqrt{a_{\max} \cdot R_{\min}} = \sqrt{0{,}735 \cdot 0{,}12} = \sqrt{0{,}0882} \approx \mathbf{0{,}297\text{ m/s}} \quad (29{,}7\text{ cm/s})
\end{equation}

\subsection{\texorpdfstring{Profil de Freinage Ciblé et Arrêt de Précision ($\pm 30\text{ mm}$)}{Profil de Freinage Ciblé et Arrêt de Précision (+/- 30 mm)}}
Pour satisfaire la spécification \textbf{S1} (immobilisation devant un obstacle à la distance cible $d_{\text{cible}} = 200\text{ mm}$ avec tolérance stricte $\pm 30\text{ mm}$), la vitesse résiduelle en fonction de la distance mesurée $d$ est régie par la loi de décélération quadratique optimale~:
\begin{equation}
    v_{\text{freinage}}(d) = \sqrt{2 \cdot a_{\max} \cdot (d - d_{\text{cible}})}
\end{equation}
Lorsque $|d - d_{\text{cible}}| \le 2\text{ cm}$, la commande commute sur l'état \code{Stopped} sans choc mécanique, garantissant la parfaite stabilité de la bille.
